% Experiment 12: Design and Implementation of RSA Public-Key Cryptosystem
\section{Theory}
RSA is an asymmetric cryptosystem based on the computational difficulty of
factoring the product of two large primes. Given distinct primes $p$ and $q$,
let $n = p \times q$ and $\phi(n) = (p-1)(q-1)$. A public exponent $e$ is
chosen with $1 < e < \phi(n)$ and $\gcd(e, \phi(n)) = 1$, and the private
exponent $d$ satisfies $d \cdot e \equiv 1 \pmod{\phi(n)}$. The public key
is $(e, n)$ and the private key is $(d, n)$. Encryption and decryption use
modular exponentiation:
\[
C = M^e \bmod n, \qquad M = C^d \bmod n
\]
RSA's security relies on the fact that, while multiplying two large primes
is computationally easy, factoring their product back into $p$ and $q$ is
computationally infeasible for sufficiently large key sizes.

\section{Algorithm}
\begin{enumerate}
  \item Choose two distinct large prime numbers $p$ and $q$.
  \item Compute $n = p \times q$.
  \item Compute $\phi(n) = (p-1)(q-1)$.
  \item Choose public exponent $e$ such that $1 < e < \phi(n)$ and $\gcd(e, \phi(n)) = 1$.
  \item Compute private exponent $d$ as the modular inverse of $e$ modulo $\phi(n)$.
  \item Publish the public key $(e, n)$ and keep the private key $(d, n)$ secret.
  \item To encrypt message $M$: compute $C = M^e \bmod n$.
  \item To decrypt ciphertext $C$: compute $M = C^d \bmod n$.
  \item Verify that the decrypted message matches the original plaintext.
\end{enumerate}

\section{Code}
\begin{minted}{cpp}
#include<bits/stdc++.h>
using namespace std;
class rsa {
  private:
    long long p, q, n, phi, e, d;
    long long gcd(long long a, long long b) {
      return b == 0 ? a : gcd(b, a % b);
    }
    long long modPow(long long base, long long exp, long long mod) {
      long long result = 1;
      base %= mod;
      while (exp > 0) {
        if (exp & 1) result = (__int128)result * base % mod;
        base = (__int128)base * base % mod;
        exp >>= 1;
      }
      return result;
    }
    long long modInverse(long long a, long long m) {
      long long x0 = 1, x1 = 0, b = m;
      while (b != 0) {
        long long qt = a / b;
        long long t = b;
        b = a % b;
        a = t;
        t = x1;
        x1 = x0 - qt * x1;
        x0 = t;
      }
      return (x0 % m + m) % m;
    }
    bool isPrime(long long num) {
      if (num < 2) return false;
      for (long long i = 2; i * i <= num; i++)
        if (num % i == 0) return false;
      return true;
    }
  public:
    rsa(long long primeP, long long primeQ);
    long long encrypt(long long msg);
    long long decrypt(long long cipher);
};
rsa::rsa(long long primeP, long long primeQ) {
  p = primeP;
  q = primeQ;
  while (!isPrime(p) || !isPrime(q) || p == q) {
    cout << "p and q must be distinct primes. Enter again: ";
    cin >> p >> q;
  }
  n = p * q;
  phi = (p - 1) * (q - 1);
  e = 2;
  while (e < phi && gcd(e, phi) != 1) e++;
  d = modInverse(e, phi);
  cout << "n   = " << n << endl;
  cout << "phi = " << phi << endl;
  cout << "Public Key  (e, n) = (" << e << ", " << n << ")" << endl;
  cout << "Private Key (d, n) = (" << d << ", " << n << ")" << endl;
}
long long rsa::encrypt(long long msg) {
  return modPow(msg, e, n);
}
long long rsa::decrypt(long long cipher) {
  return modPow(cipher, d, n);
}
int main() {
  cout << "Enter two prime numbers p and q: ";
  long long p, q;
  cin >> p >> q;
  rsa cipher(p, q);

  cout << "\nEnter plaintext message (integer < n): ";
  long long msg;
  cin >> msg;

  long long cipherText = cipher.encrypt(msg);
  cout << "Ciphertext = M^e mod n = " << cipherText << endl;

  long long decrypted = cipher.decrypt(cipherText);
  cout << "Decrypted  = C^d mod n = " << decrypted << endl;
  return 0;
}
\end{minted}

\section{Result - I/O}
\begin{verbatim}
Enter two prime numbers p and q: 23 97
n   = 2231
phi = 2112
Public Key  (e, n) = (5, 2231)
Private Key (d, n) = (845, 2231)

Enter plaintext message (integer < n): 1024

--- Encryption ---
Ciphertext = M^e mod n = 1168

--- Decryption ---
Decrypted  = C^d mod n = 1024
\end{verbatim}

\section{Conclusion}
The RSA cryptosystem was successfully implemented, covering key generation,
encryption, and decryption via modular exponentiation and the extended Euclidean
algorithm. The experiment validates RSA's core property that
$(M^e)^d \equiv M \pmod{n}$. This implementation uses small values for
demonstration; production RSA uses much larger primes and padding schemes.
